\section{Contexto}
\label{sec:contexto}

A computação tem passado a integrar os currículos escolares de diferentes países como uma área de conhecimento, e não apenas como um conjunto de ferramentas utilizadas em outras disciplinas \cite{vegas2021building}. Seu ensino envolve conceitos, práticas e formas de raciocínio relacionados, entre outros temas, à resolução de problemas, à representação e ao tratamento de informações, ao funcionamento dos sistemas digitais e à compreensão dos impactos sociais das tecnologias.

No Brasil, esse movimento ganhou forma normativa em 2022, com a instituição de normas de Computação, em complemento à \sigla{BNCC}{Base Nacional Comum Curricular}. A BNCC define as aprendizagens essenciais da Educação Básica, enquanto o complemento estabelece competências, objetos de conhecimento, objetivos de aprendizagem e habilidades de computação para a educação infantil, o ensino fundamental e o ensino médio \cite{brasil2022resolucao1,brasil2022bncccomputacao}. Neste trabalho, a expressão BNCC Computação é utilizada para se referir a esse conjunto normativo.

A criação dessas aprendizagens gera uma demanda nacional de implementação. Estados, municípios, redes de ensino e escolas precisam decidir como a computação será incorporada aos seus currículos e às suas práticas pedagógicas em todos os níveis escolares. Isso envolve definir como os conhecimentos serão organizados, quem será responsável por ensiná-los e quais condições serão necessárias para que o ensino ocorra.

A formação de professores é uma dessas condições. A Resolução CNE/CEB nº 1/2022 determina que a formação inicial e a formação continuada considerem as aprendizagens previstas na BNCC Computação \cite{brasil2022resolucao1}. Entretanto, essa determinação não especifica como os conhecimentos da área devem ser incorporados aos currículos das licenciaturas. Também não indica quais cursos devem incorporá-los, quais conhecimentos são pertinentes a cada perfil profissional, com que profundidade devem ser abordados ou em quais espaços curriculares (disciplinas específicas, componentes já existentes, práticas ou estágios) devem ser inseridos.

Esta pesquisa situa-se nesse ponto de encontro entre formação inicial de professores, que pode ser compreendida como o processo desenvolvido ao longo de todo o curso de licenciatura para preparar o estudante para o exercício da docência, e a implementação da BNCC Computação. Seu foco está nas decisões necessárias para integrar conhecimentos de Computação a licenciaturas já existentes, considerando tanto as finalidades formativas de cada curso quanto suas condições curriculares e institucionais.


\section{Motivação}
\label{sec:motivacao}

A existência de uma norma curricular não garante que as aprendizagens nela previstas sejam efetivamente ensinadas, pois o currículo desenvolvido nas escolas pode diferir daquele estabelecido nos documentos oficiais \cite{remillard2014conceptualizing}. A passagem do currículo prescrito para a prática depende de condições relacionadas à organização curricular, aos conhecimentos específicos e pedagógicos dos professores, ao tempo disponível para planejamento e formação, aos materiais e recursos e ao apoio oferecido pelas instituições \cite{yiatas2026enacting,royal2017afterreboot}. Essas condições não atuam isoladamente. A formação docente, por exemplo, precisa ser acompanhada de oportunidades de planejamento, apoio institucional e condições adequadas de trabalho. Da mesma forma, a simples disponibilidade de equipamentos, plataformas ou materiais não assegura a implementação do ensino de Computação quando os professores não dispõem de clareza curricular, preparação pedagógica e critérios para selecionar e utilizar esses recursos \cite{royal2017afterreboot}.

Um estudo preliminar, coautorado pelo candidato, identificou esse descompasso nas redes públicas paulistas \cite{santana2026desafios}. Os resultados indicaram que a dificuldade não se limita ao conhecimento da norma. O problema principal está em traduzi-la para o currículo e para a rotina escolar, diante de uma formação fragmentada, de uma infraestrutura desigual e de responsabilidades institucionais pouco definidas. O estudo também mostrou o risco de interpretar a implementação da BNCC Computação apenas como uma ampliação do uso de plataformas ou de recursos digitais, sem tratar a computação como objeto de aprendizagem. O estudo é apresentado em mais detalhes na Seção~\ref{sec:estudo-implementacao-bncc}.

Para os professores que já atuam na educação básica, a formação continuada constitui a resposta mais imediata a essa demanda. Ela é necessária para apresentar novos conteúdos, apoiar mudanças curriculares e acompanhar as transformações da área. Entretanto, parte das iniciativas documentadas na literatura ocorre por meio de cursos ou oficinas de curta duração e, recentemente, por cursos de pós-graduação como o \sigla{PROFCOMP}{Mestrado Profissional em Ensino de Computação} \footnote{\url{https://www.sbc.org.br/profcomp-sbc/}}. Embora essas formações possam produzir ganhos de conhecimento e de confiança, sua transformação em práticas permanentes depende do apoio recebido posteriormente nas escolas e redes de ensino \cite{liu2024bringing,ma2025effectiveness,mouza2022investigating}.

A formação continuada, portanto, não deve ser tratada como a única resposta. Quando os conhecimentos necessários não foram desenvolvidos na formação inicial, cursos pontuais precisam compensar as lacunas conceituais e pedagógicas acumuladas ao longo da trajetória profissional. A formação inicial oferece condições diferentes. Por estar distribuída ao longo de vários semestres, pode permitir a continuidade, a progressão e a articulação entre conhecimentos de computação, conhecimentos pedagógicos, conhecimentos específicos (ex. matemática, física, artes) e experiências práticas em sala de aula.

Essa preparação não precisa ser idêntica em todas as licenciaturas. A implementação da BNCC Computação pode envolver professores especialistas, professores polivalentes e docentes de diferentes componentes curriculares. Esses profissionais atuam em etapas distintas da Educação Básica, possuem formações disciplinares diferentes e podem assumir responsabilidades também diferentes. Um professor formado em Pedagogia, por exemplo, costuma atuar de forma polivalente, responsável por todos os componentes curriculares de uma turma na educação infantil ou nos anos iniciais do ensino fundamental, enquanto um licenciado em matemática, física, química ou computação é formado especificamente para ensinar essa área como disciplina própria nos anos finais do ensino fundamental e no ensino médio. Assim, para diferentes licenciaturas, é necessário decidir quais conhecimentos de computação são relevantes, com que profundidade devem ser trabalhados, em que momento ou espaço do currículo devem ser inseridos e qual carga horária pode ser destinada a essa formação, a fim de atender às demandas da BNCC Computação.

Transformar essa necessidade em mudanças nos cursos de formação inicial não é uma tarefa simples. As licenciaturas (fora da área da computação) possuem currículos em funcionamento, cargas horárias definidas, componentes obrigatórios, perfis de egresso, regras de aprovação e corpos docentes com formações específicas. Uma nova demanda precisa disputar espaço com outras necessidades formativas e ser analisada por coordenações, Núcleos Docentes Estruturantes, colegiados e demais instâncias institucionais. A alternativa mais abrangente nem sempre é a mais viável, enquanto uma alteração de baixa complexidade pode não produzir formação suficiente.

Faz-se necessário, portanto, compreender e apoiar a forma como as equipes curriculares transformam essa demanda geral em decisões concretas sobre os conhecimentos a serem trabalhados, sua profundidade, os momentos ou espaços curriculares e a carga horária, considerando as condições específicas de cada curso.

\section{Problema de pesquisa}
\label{sec:problema-pesquisa}

A BNCC Computação estabelece aprendizagens para toda a educação básica e determina que a formação de professores as considere. As normas, entretanto, não definem quais licenciaturas devem incorporar conhecimentos de computação, quais conhecimentos são relevantes para cada uma, com que profundidade devem ser trabalhados, em que momento ou espaço curricular devem ser inseridos, com qual carga horária ou sob responsabilidade de quais profissionais. Também não indicam como essas decisões devem considerar o currículo vigente, os recursos disponíveis e as restrições de cada instituição.

Esse problema é agravado pela situação atual dos cursos. Fora da Licenciatura em Computação, a presença de conhecimentos próprios da área tende a ser limitada, optativa, descontínua ou associada ao uso instrumental de tecnologias \cite{delgado2026computacao}. Ao mesmo tempo, incluir novos conhecimentos em uma licenciatura de outra área exige decisões interdependentes sobre objetivos formativos, profundidade, progressão, espaços curriculares, carga horária, práticas pedagógicas, recursos e viabilidade institucional. A delimitação dessa lacuna é desenvolvida nos capítulos ~\ref{chapter:referencial-teorico} e~\ref{chapter:revisao-integracao-curricular}.

A literatura sobre educação em computação já diagnostica a ausência da Computação nos currículos \cite{delgado2026computacao}, propõe referenciais conceituais sobre o que ensinar e desenvolve modelos curriculares para orientar sua inserção na Educação Básica \cite{k2016k, NewZealandMoE2017, ribeiro2019diretrizes, abrantes2024pensamento}. Contudo, ainda são escassos os trabalhos que orientam a transformação desses referenciais gerais, como a BNCC Computação, em decisões curriculares contextualizadas para cursos de licenciatura que não têm a Computação como área principal de formação.

Dessa forma, torna-se importante criar estratégias que apoiem o processo pelo qual uma equipe curricular analisa as necessidades de uma licenciatura específica, articula referenciais normativos e curriculares, incluindo computação, considera o currículo vigente e as condições institucionais e transforma uma necessidade geral de integração da computação em decisões concretas e justificadas.

\subsection*{Objetivos da pesquisa:}

O objetivo geral da pesquisa é:

\begin{quote}
Desenvolver e avaliar um protocolo sistematizado, implementado em uma ferramenta computacional, que apoie equipes curriculares de licenciaturas que não possuem a Computação como área principal de formação na formulação de decisões contextualizadas e justificadas sobre a integração da Computação, considerando os conhecimentos relevantes, sua profundidade, os espaços curriculares, a carga horária e as condições específicas de cada curso.
\end{quote}

Para alcançar esse objetivo, foram definidos os seguintes objetivos específicos:

\begin{enumerate}

\item Caracterizar o problema de integração curricular da Computação em licenciaturas brasileiras, considerando a literatura, as diretrizes normativas, as dificuldades de implementação e os conhecimentos formativos relevantes;
\item Mapear como licenciaturas selecionadas contemplam atualmente conhecimentos relacionados à Computação, considerando conteúdos, profundidade, obrigatoriedade, carga horária, posição no curso e articulação entre componentes curriculares;
\item Definir os objetivos e os requisitos da solução, relacionando as necessidades identificadas às funções, aos critérios, às informações de entrada e aos resultados esperados do protocolo;
\item Projetar, desenvolver e implementar a primeira versão do protocolo e da ferramenta computacional de apoio à sua aplicação;
\item Demonstrar e avaliar o protocolo em situações representativas e em licenciaturas reais que não possuem a Computação como área principal de formação, examinando sua clareza, consistência, utilidade, aplicabilidade, adequação ao contexto e capacidade de produzir recomendações justificadas;
\item Refinar e consolidar o protocolo com base nas evidências obtidas, documentando suas versões, limitações, condições de aplicação e contribuições práticas e científicas.
\end{enumerate}

\section{Organização do trabalho}
\label{sec:organizacao-trabalho}

Além deste capítulo introdutório, este trabalho está organizado da seguinte forma. O Capítulo~\ref{chapter:referencial-teorico} apresenta os fundamentos conceituais, normativos e educacionais da pesquisa. Inicialmente, delimita a Computação como um campo de conhecimentos e de aprendizagem na Educação Básica. Em seguida, examina a BNCC Computação, as responsabilidades e desafios decorrentes de sua implementação. O capítulo também discute o papel da formação docente, os limites de uma resposta baseada exclusivamente na formação continuada e as possibilidades oferecidas pela formação inicial.

O Capítulo~\ref{chapter:revisao-integracao-curricular} concentra a discussão sobre a integração da Computação aos currículos das licenciaturas. O capítulo examina as diferenças entre os perfis profissionais formados por esses cursos. Em seguida, analisa a integração da Computação em currículos como um problema de decisão curricular. Por fim, são analisadas a presença atual da Computação nas licenciaturas e as contribuições oferecidas pela literatura, o que permite delimitar a lacuna que fundamenta a pesquisa.

O Capítulo~\ref{chapter:proposta} apresenta a solução proposta para enfrentar essa lacuna: um protocolo sistematizado, implementado em uma ferramenta computacional, para apoiar decisões sobre a integração da Computação em cursos de licenciatura. O capítulo define as perguntas e os objetivos da pesquisa, descreve o artefato pretendido e apresenta o delineamento metodológico baseado em \textit{Design Science Research}. Também são detalhadas as atividades de construção, demonstração, avaliação e refinamento do protocolo, bem como os aspectos éticos, as limitações previstas, as ameaças à validade, o cronograma e os resultados esperados.

Por fim, o Capítulo sistematiza as atividades acadêmicas e formativas já desenvolvidas em temas relacionados ao ensino de Computação e à implementação da BNCC Computação. O capítulo apresenta os artigos publicados no decorrer da pesquisa, cursos de formação ministrados a docentes e as disciplinas cursadas na área de ensino de Computação.